\section*{Conclusion du chapitre}
\addcontentsline{toc}{section}{Conclusion du chapitre}

Ce chapitre avait pour objectif de présenter la construction du modèle thématique et de montrer de quelle manière ses résultats peuvent contribuer à l'étude du naturalisme dans l'œuvre romanesque d'Émile Zola. L'utilisation de BERTopic a permis d'associer une représentation sémantique contextualisée des segments textuels à des méthodes de réduction dimensionnelle, de regroupement et de caractérisation lexicale. La sélection du modèle ne s'est ainsi pas limitée à la recherche d'une configuration produisant un nombre déterminé de topics : elle s'est appuyée sur l'évaluation de 256 configurations, sur un classement multicritère et sur une vérification de la stabilité des résultats obtenus avec plusieurs graines aléatoires.

Le choix final d'un modèle composé de 33 topics répond à la volonté de préserver la diversité thématique du corpus. Une partition réduite à huit ensembles aurait certes facilité la lecture générale, mais elle aurait également regroupé des réalités sociales, matérielles et narratives distinctes. Les 33 topics retenus permettent au contraire de différencier plus précisément les espaces sociaux, les métiers, les institutions, les rapports économiques, les objets du quotidien ainsi que les formes de souffrance représentées dans les romans. La validation mathématique et l'examen de la stabilité ne garantissent pas une vérité définitive du découpage, mais ils montrent que celui-ci ne résulte pas uniquement d'une configuration accidentelle de l'algorithme.

L'interprétation des dix termes les plus représentatifs de chaque topic a ensuite permis d'établir plusieurs degrés de proximité avec le naturalisme. Cette gradation évite de réduire celui-ci à une opposition binaire entre topics naturalistes et non naturalistes. Certains ensembles lexicaux renvoient directement à des composantes centrales de l'esthétique naturaliste, notamment le travail, la condition populaire, l'argent, les déterminismes sociaux ou corporels et les formes de souffrance. D'autres entretiennent un rapport plus indirect avec ce cadre esthétique, tandis que certains topics relèvent surtout de situations narratives particulières. Le naturalisme apparaît donc moins comme un thème isolé que comme une logique transversale reliant plusieurs dimensions de l'univers zolien.

Les visualisations temporelles complètent cette lecture en replaçant les topics dans la chronologie de la carrière de Zola. Elles montrent que les principales préoccupations associées au naturalisme ne sont pas cantonnées à un seul roman ni à une phase unique de son œuvre. Leur fréquence varie selon les périodes et selon les objets privilégiés par chaque récit, mais elles réapparaissent sous différentes formes tout au long du corpus. Ces variations invitent à distinguer la permanence d'une préoccupation de l'uniformité de son expression : un topic peut connaître des phases de forte concentration, s'atténuer, puis réapparaître dans un contexte narratif différent.

Cette analyse confirme ainsi l'intérêt de la modélisation thématique pour étudier un corpus littéraire étendu sans substituer le calcul à l'interprétation. BERTopic permet de repérer des régularités, des proximités et des évolutions difficilement observables à l'échelle de l'ensemble des romans, tandis que leur signification demeure dépendante d'une lecture critique des mots représentatifs et des textes concernés. La suite du mémoire pourra dès lors approfondir ces résultats en confrontant les tendances quantitatives aux choix esthétiques, sociaux et narratifs propres à l'écriture zolienne.